\begin{mistake}{2026-06-30}{03}

\begin{problemcontent}
设函数 \(f(x)\) 在 \([0,1]\) 上有连续一阶导数，且 \(f(0)=0\)。试证明：至少存在一点 \(\xi\in[0,1]\)，使得
\[
f'(\xi)=2\int_0^1 f(x)\,dx.
\]
\end{problemcontent}

\begin{solutioncontent}
由 \(f(0)=0\) 及牛顿--莱布尼茨公式，有
\[
f(x)=f(x)-f(0)=\int_0^x f'(t)\,dt.
\]
于是
\[
\int_0^1 f(x)\,dx
=\int_0^1\left(\int_0^x f'(t)\,dt\right)dx.
\]
对积分区域 \(0\le t\le x\le 1\) 交换积分次序，得
\[
\int_0^1\left(\int_0^x f'(t)\,dt\right)dx
=\int_0^1\left(\int_t^1 f'(t)\,dx\right)dt
=\int_0^1(1-t)f'(t)\,dt.
\]
因为 \(f'\) 在 \([0,1]\) 上连续，且 \(1-t\ge 0\)，由加权积分中值定理，存在 \(\xi\in[0,1]\)，使
\[
\int_0^1(1-t)f'(t)\,dt
=f'(\xi)\int_0^1(1-t)\,dt.
\]
又
\[
\int_0^1(1-t)\,dt=\frac12,
\]
所以
\[
\int_0^1 f(x)\,dx=\frac12 f'(\xi),
\]
即
\[
f'(\xi)=2\int_0^1 f(x)\,dx.
\]
故命题得证。
\end{solutioncontent}

\mistakeknowledge{
\begin{itemize}
  \item \textbf{核心公式/定理}：由 \(f(0)=0\) 得 \(f(x)=\int_0^x f'(t)\,dt\)；加权积分中值定理：若 \(\varphi\) 连续、\(g\ge0\)，则 \(\int_a^b\varphi g=\varphi(\xi)\int_a^b g\)。
  \item \textbf{易错点}：交换积分次序时，区域 \(0\le t\le x\le1\) 应改写为 \(0\le t\le1,\ t\le x\le1\)，从而得到权函数 \(1-t\)。
  \item \textbf{关联知识}：本题也可用构造函数结合罗尔定理，或用拉格朗日中值定理与导数最大最小值夹逼证明。
\end{itemize}}

\end{mistake}
